\section{Composition and differentiation in a reciprocal logarithm}
\label{sec:reciprocal-log-operations}

The reciprocal-logarithmic unit of Section~\ref{sec:logarithmic-scales}
has more regularity than an arbitrary function satisfying a magnitude
Big-$O$ estimate. For an exact rational logarithmic forward unit, its
implicit inverse is holomorphic in the small reciprocal logarithm. This
regularity supplies composition and differentiation rules without
discarding the exact power carrier or asking for an unsupported derivative
of an unknown remainder.

Consider exact rational logarithmic forward units with positive monomial
inverse carriers and zero finite source and target offsets. Known positive
source-power perturbations are not included in this particular analytic
construction: although smaller than every reciprocal-logarithmic order,
they require a separate differentiable-remainder argument. The composition
and differentiation below preserve holomorphy of the coefficient function.

\subsection{A common positive logarithmic coordinate}

Write
\begin{equation}
 t=\frac{\epsilon}{\log y}>0,\qquad
 \epsilon=\begin{cases}-1,&y\to0^+,\\1,&y\to+\infty.\end{cases}
 \qquad g(y)=C y^\alpha A(t),
 \label{eq:reciprocal-calculus-carrier}
\end{equation}
Here $C>0$, $\alpha$ is a fixed real number, and $A$ is holomorphic
near $t=0$. For an original inverse unit, $A(0)=1$. Differentiation can
produce a coefficient function whose constant vanishes or has either sign;
these outcomes remain valid carriers for further differentiation and as
outer functions in composition.

Suppose the original inverse has positive leading amplitude $a$ and
internal leading source power $p\ne0$. Its small coordinate is
$\tau=-1/\log((y/a)^{1/p})$, and therefore
\begin{equation}
 \tau=\frac{|p|t}{1-\epsilon(\log a)t},\qquad
 \epsilon=-\operatorname{sign}p.
 \label{eq:reciprocal-calculus-rechart}
\end{equation}
This is an analytic change of variable at zero with positive derivative.
It transports an analytic $\OO(\tau^\beta)$ tail to
$\OO(t^\beta)$. For an internal observable power $r_*$, the monomial
carrier has $\alpha=r_*/p$ and $C=\sigma^r a^{-\alpha}$.
Assume this $C$ is positive. In addition to the
original target domain, the canonical chart requires $y>0$ and
$\epsilon\log y>0$. These are eventual branch conditions; they are not
a computed radius of convergence.

The Taylor expansions have nonnegative integer powers of $t$ and constant real
coefficients. A finite exclusive cutoff $h$ retains powers $n<h$; the
error exponent is limited by the inherited Taylor remainder and any
discarded coefficient. Cancellation can make the first omitted nonzero
degree larger than the cutoff.

\subsection{Composition of two monomial carriers}

Let the outer function and inner function be
\[
 g(v)=C_o v^\alpha A\!\left(\frac{\epsilon_o}{\log v}\right),
 \qquad h(y)=C_i y^b V(t),\qquad
 t=\frac{\epsilon_i}{\log y},\quad V(0)=1.
\]
A positive nonunit constant in the inner coefficient function is first
absorbed into $C_i$. The composition requires $b\ne0$ and
\begin{equation}
 \epsilon_o b\epsilon_i>0.
 \label{eq:reciprocal-composition-endpoints}
\end{equation}
This condition says that the positive inner function approaches the outer
function's selected endpoint. It permits a negative power to interchange
zero and infinity when the two specified endpoints agree with that change.

Taking the real logarithm of the inner carrier gives the exact identities
\begin{align}
 T(t)&=\frac{\epsilon_o t}
 {b\epsilon_i+t(\log C_i+\log V(t))},
 \label{eq:reciprocal-composed-coordinate}\\
 g(h(y))&=C_o C_i^\alpha y^{\alpha b}
       V(t)^\alpha A(T(t)).
 \label{eq:reciprocal-composition-identity}
\end{align}
The denominator at zero is nonzero, $T(0)=0$, and
$T'(0)=\epsilon_o/(b\epsilon_i)>0$. The functions $\log V$ and
$V^\alpha$ use their analytic branches with $V(0)=1$. Consequently the
new coefficient function is holomorphic near zero and the original
monomial carrier remains exact.

If the outer and inner Taylor remainders are
$\OO(T^{\beta_o})$ and $\OO(t^{\beta_i})$, respectively, a conservative
relative remainder for \eqref{eq:reciprocal-composition-identity} is
\begin{equation}
 \OO\!\left(t^{\min(\beta_o,\beta_i)}\right).
 \label{eq:reciprocal-composition-error}
\end{equation}
To see this, an inner perturbation of size $\OO(t^{\beta_i})$ changes
$\log V$ and $V^\alpha$ by that order, while it changes $T$ by
$\OO(t^{\beta_i+2})$. The outer function has bounded derivatives near
zero, and its original remainder is evaluated at a coordinate comparable
to $t$. The Taylor algebra preserves these bounds and can improve them when
a constant factor or derivative vanishes.
Truncating the displayed polynomial adds its own first discarded order.
A larger formal cutoff never removes either operand's unknown tail.

For example, if both functions have the form $yW(-1/\log y)$ near zero,
the coefficient function of their composition is
\[
 W(t)\,W\!\left(\frac{t}{1-t\log W(t)}\right).
\]
At positive infinity the denominator is instead $1+t\log W(t)$.
These formulas determine the coefficients by ordinary Taylor expansion.
For general powers and scales, \eqref{eq:reciprocal-composed-coordinate}
retains both the exponent factor $b$ and the additive logarithm
$\log C_i$; dropping either gives incorrect coefficients already at low
reciprocal-logarithmic orders.

An inner coefficient function with a zero constant term is excluded from
this composition rule. Its logarithm introduces an additional $\log t$
scale. A negative constant also fails the positive inner branch condition.
Those failures describe a change in the required scale or branch, rather
than a reason to take a logarithm of an unproved positive approximation.

\subsection{Differentiating the analytic remainder}

The coordinate derivative is
\[
 \frac{dt}{dy}=-\frac{\epsilon t^2}{y}.
\]
Hence
\begin{equation}
 \frac{d}{dy}\bigl(Cy^\alpha A(t)\bigr)
 =Cy^{\alpha-1}\left(\alpha A(t)-\epsilon t^2 A'(t)\right).
 \label{eq:reciprocal-carrier-derivative}
\end{equation}
For the $n$th derivative, the coefficient function is obtained by applying
the successive operators
\begin{equation}
 \prod_{k=0}^{n-1}
 \left(\alpha-k-\epsilon t^2\frac{d}{dt}\right)
 \quad\hbox{to }A(t),
 \label{eq:reciprocal-carrier-higher-derivatives}
\end{equation}
and the carrier is $Cy^{\alpha-n}$. The factors differ only by constants
and commute. Applying them successively displays the precision change
and possible cancellation at each derivative order.

If a holomorphic remainder has a zero of order at least $\beta$ at zero,
then its $j$th derivative is $\OO(t^{\beta-j})$ for every fixed $j$.
This follows directly from its convergent Taylor series, or from Cauchy
estimates on a smaller disk. Thus the $t^2 A'$ contribution in
\eqref{eq:reciprocal-carrier-derivative} has remainder
$\OO(t^{\beta+1})$, while a nonzero $\alpha A$ contribution has remainder
$\OO(t^\beta)$. If $\alpha=0$, the latter contribution disappears and the
relative logarithmic error improves by one power. This is a proved
derivative contract for the exact rational-logarithmic implicit model;
it is not a rule for differentiating an arbitrary magnitude Big-$O$.

For the inverse of $x+x/\log x$ near zero, the retained coefficient
function begins
\[
 A(t)=1+t+t^2+2t^3+\tfrac72t^4+\OO(t^5).
\]
The first derivative therefore has the expansion
\[
 1+t+2t^2+4t^3+\tfrac{19}{2}t^4+\OO(t^5).
\]
Its carrier is constant. Differentiating once more gives a carrier $y^{-1}$
and a relative remainder $\OO(t^6)$, provided the known terms generated
through degree five are retained. This illustrates the precision improvement
when a carrier derivative vanishes. At infinity the signs of the odd
powers and the coordinate derivative change consistently.

\subsection{Refinement and the analytic hypothesis}

A stronger approximation to a composed or differentiated function requires
additional coefficients of its underlying analytic functions. Merely
truncating the same finite polynomial at a larger formal degree supplies
no new information about its unknown tail. Once those extra coefficients
are known, the exact substitution formulas and derivative operators above
transport their precision to the result.

The remainder theorem depends on holomorphy in the reciprocal logarithm,
rather than on finite coefficient identities alone. A source-power sector
which is smaller than every power of $t$ can be invisible to the Taylor
series yet obstruct holomorphy at $t=0$. Its differentiation therefore
requires a separate sector estimate. The hypotheses identifying an exact
rational logarithmic model are part of the mathematical statement and
cannot be discarded after its finite Taylor coefficients have been found.

As with the underlying implicit-function theorem, the convergence radius
and bounding constants exist for fixed data. The formulas in this section
do not assign explicit numerical values to those constants.
